\chapter{Module-LWE (MLWE)}
\label{ch:mlwe}

\section{The three generations of lattice-based cryptography}

We have already seen three major stages in the development of lattice-based cryptography.

\begin{itemize}
    \item \textbf{First generation: LWE}. The equation is
    \[
    b = A s + e,
    \]
    where $A$ is a random matrix. The security proof is elegant, but the representation is large and the computation is costly.

    \item \textbf{Second generation: RLWE}. The matrix is replaced by a polynomial ring structure. Computation becomes much faster, but the algebraic structure becomes stronger and more specialized.

    \item \textbf{Third generation: Module-LWE}. This is a compromise between the large but simple LWE construction and the compact but highly structured RLWE construction.
\end{itemize}

The historical path is therefore:

\begin{center}
\begin{tikzpicture}[node distance=9mm, every node/.style={font=\small}]
  \node[flownode] (a) {LWE};
  \node[flownode, right=of a] (b) {RLWE};
  \node[keynode, right=of b] (c) {Module-LWE};
  \draw[flowarrow] (a) -- (b);
  \draw[flowarrow] (b) -- (c);
\end{tikzpicture}
\end{center}

\section{What is a module?}

The word ``module'' is not informal jargon; it is a precise algebraic structure, and Kyber's secrets live in one. It generalizes the vector space of Chapter~\ref{ch:linalg} by letting the scalars come from a ring instead of a field.

\begin{definition}[Module over a ring]
Let $R$ be a commutative ring with identity $1$. A (left) \emph{$R$-module} is an abelian group $(M,+)$ together with a scalar multiplication $R\times M\to M$, written $(r,m)\mapsto rm$, such that for all $r,s\in R$ and $x,y\in M$:
\begin{enumerate}
  \item $r(x+y)=rx+ry$;
  \item $(r+s)x=rx+sx$;
  \item $(rs)x=r(sx)$;
  \item $1\,x=x$.
\end{enumerate}
\end{definition}

These are exactly the vector-space axioms with the field replaced by a ring. When $R$ is a field, an $R$-module \emph{is} a vector space; a module is the strictly more general object obtained by relaxing ``field'' to ``ring.'' Several properties carry over, and one convenience is lost.

\begin{itemize}
  \item \textbf{Free modules and rank.} The set $R^{k}=\{(m_1,\dots,m_k):m_i\in R\}$ with componentwise operations is an $R$-module, the \emph{free module of rank $k$}. Like a vector space it has a basis -- the standard unit vectors $e_1,\dots,e_k$ -- so every element is a unique $R$-combination of them. Kyber's secret space is exactly the free module $R_q^{k}$.
  \item \textbf{Matrices act as homomorphisms.} A matrix $A\in R^{k\times k}$ defines an $R$-linear map $R^{k}\to R^{k}$, $\mathbf{x}\mapsto A\mathbf{x}$, called an \emph{$R$-module homomorphism}. The Module-LWE map $\mathbf{s}\mapsto A\mathbf{s}$ is one of these.
  \item \textbf{Submodules and ideals.} A submodule is the analogue of a subspace; the submodules of $R$ regarded as a module over itself are precisely its \emph{ideals} (Chapter~\ref{ch:polyring}) -- which is why RLWE is said to rest on \emph{ideal} lattices.
  \item \textbf{What is lost.} Over a general ring, not every module is free and not every submodule has a basis -- a real departure from vector spaces. Kyber avoids this entirely by working in the free module $R_q^{k}$.
\end{itemize}

With the structure named precisely, the three-generation progression is a statement about \emph{which module the secret lives in}:
\[
\underbrace{s\in\mathbb{Z}_q^{n}}_{\text{LWE}}
\qquad
\underbrace{s(x)\in R_q}_{\text{RLWE}}
\qquad
\underbrace{\mathbf{s}\in R_q^{k}}_{\text{Module-LWE}}.
\]
LWE uses a vector of field elements; RLWE uses a single ring element; Module-LWE uses a length-$k$ vector over the ring $R_q$ -- a free module of rank $k$.

\section{Kyber's secret}

In ML-KEM (Kyber), the secret is not a single polynomial but a vector of polynomials. The security level determines the dimension of that vector.

\begin{itemize}
    \item Kyber512 uses $k=2$.
    \item Kyber768 uses $k=3$.
    \item Kyber1024 uses $k=4$.
\end{itemize}

Thus, the security level is tied to the size of the module.

\section{Why does the matrix come back?}

At first glance, people are surprised to see matrices again in Kyber. But the matrix is not a large random matrix over integers. Instead, its entries are polynomials.

For example, for $k=2$, we may have
\[
A =
\begin{bmatrix}
 a_{11}(x) & a_{12}(x) \\
 a_{21}(x) & a_{22}(x)
\end{bmatrix}.
\]

Each entry is a polynomial in the ring $R_q$. The object is still a matrix, but the matrix entries are now ring elements rather than ordinary numbers.

\section{Kyber public key generation}

The public key generation is very similar to LWE in spirit. We choose a secret vector $\mathbf{s}$, a noise vector $\mathbf{e}$, and a polynomial matrix $A$, and then compute
\[
\mathbf{t} = A\mathbf{s} + \mathbf{e}.
\]

This is exactly the same governing idea as LWE, except that the scalars are now ring elements. The figure below makes the shapes of $A$ and $\mathbf{s}$ concrete for the smallest case $k=2$ (indices run from $0$ to $k-1$, matching the Sage code below).

\begin{figure}[H]
\centering
\begin{tikzpicture}[every node/.style={font=\small}]
  % A 2x2 matrix of ring elements: each entry a_{ij}(x) is itself a
  % polynomial in R_q, not a single scalar as it would be for plain LWE
  \node[draw=pqcblue, thick, fill=pqclight, minimum width=2.0cm, minimum height=1.0cm] (a00) at (0,1.1) {$a_{00}(x)$};
  \node[draw=pqcblue, thick, fill=pqclight, minimum width=2.0cm, minimum height=1.0cm] (a01) at (2.2,1.1) {$a_{01}(x)$};
  \node[draw=pqcblue, thick, fill=pqclight, minimum width=2.0cm, minimum height=1.0cm] (a10) at (0,-0.1) {$a_{10}(x)$};
  \node[draw=pqcblue, thick, fill=pqclight, minimum width=2.0cm, minimum height=1.0cm] (a11) at (2.2,-0.1) {$a_{11}(x)$};
  \node[fit=(a00)(a01)(a10)(a11), inner sep=6pt, draw=pqcgray, thick] (Abox) {};
  \node[pqcblue, font=\small] at ($(Abox.south)+(0,-0.5)$) {$A \in R_q^{2\times 2}$};

  % A length-2 vector of ring elements: the secret s, one entry per row of A
  \node[draw=pqcorange, thick, fill=white, minimum width=2.0cm, minimum height=1.0cm] (s0) at (5.6,1.1) {$s_0(x)$};
  \node[draw=pqcorange, thick, fill=white, minimum width=2.0cm, minimum height=1.0cm] (s1) at (5.6,-0.1) {$s_1(x)$};
  \node[fit=(s0)(s1), inner sep=6pt, draw=pqcgray, thick] (sbox) {};
  \node[pqcorange, font=\small] at ($(sbox.south)+(0,-0.5)$) {$\mathbf{s} \in R_q^{2}$};
\end{tikzpicture}
\caption{The module structure made concrete for $k=2$: the public matrix $A \in R_q^{2\times2}$ (blue, left) has four entries $a_{00}(x),a_{01}(x),a_{10}(x),a_{11}(x)$, each a full polynomial in $R_q$, and the secret $\mathbf{s}=(s_0(x),s_1(x))$ (orange, right) is a length-$2$ vector of such polynomials. Comparing this against the earlier ``Secret object'' column: plain LWE would put a single scalar in each box, and RLWE would collapse the whole picture to just one box.}
\label{fig:mlwe-module-structure}
\end{figure}

\section{Why is this better than plain LWE?}

We can compare the three designs.

\begin{itemize}
    \item \textbf{LWE}: the matrix is huge and the public key is large.
    \item \textbf{RLWE}: the representation is compact, but the algebraic structure is very strong.
    \item \textbf{Module-LWE}: the matrix is small and the elements are polynomials, giving a very good compromise between compactness and flexibility.
\end{itemize}

In other words, Module-LWE uses a small matrix of polynomials instead of a huge matrix of integers or a single large polynomial.

\section{Sage construction of a toy module}

We can create a very small toy example in Sage. First define the ring.

\begin{lstlisting}[language=Python]
q = 17
R.<x> = PolynomialRing(GF(q))
S = R.quotient(x^8 + 1)
\end{lstlisting}

Now define a secret as a vector of two polynomials.

\begin{lstlisting}[language=Python]
from random import choice
def small_poly():
    return S([choice([-1, 0, 1]) for _ in range(8)])

s = vector(S, [small_poly(), small_poly()])
print(s)
\end{lstlisting}

Create a small polynomial matrix.

\begin{lstlisting}[language=Python]
A = Matrix(S, [[S.random_element(), S.random_element()],
               [S.random_element(), S.random_element()]])
print(A)
\end{lstlisting}

Create a noise vector.

\begin{lstlisting}[language=Python]
e = vector(S, [small_poly(), small_poly()])
assert all(c in [0, 1, 16] for p in list(s) + list(e) for c in p.lift().list())
\end{lstlisting}

Then compute the public-like value.

\begin{lstlisting}[language=Python]
t = A * s + e
print(t)
\end{lstlisting}

This has the algebraic shape of a Module-LWE sample. The public matrix is uniform, while the secret and error use small coefficients; it is still a toy example rather than the FIPS 203 sampler.

\section{The Module-LWE problem}

With the module structure in place, Module-LWE is the LWE distribution of Chapter~\ref{ch:lwe} with the scalars promoted from $\mathbb{Z}_q$ to the ring $R_q$.

\begin{definition}[Module-LWE]
Fix the ring $R_q=\mathbb{Z}_q[x]/(x^n+1)$, a module rank $k$, and an error distribution $\chi$ over $R_q$ concentrated on small elements. For a secret $\mathbf{s}\in R_q^{k}$, the \emph{Module-LWE distribution} outputs
\[
(\mathbf{a},\,b),\qquad \mathbf{a}\leftarrow R_q^{k}\ \text{uniform},\quad e\leftarrow\chi,\quad b=\langle \mathbf{a},\mathbf{s}\rangle + e\in R_q.
\]
Stacking $m$ samples as the rows of $A\in R_q^{m\times k}$ gives the matrix form $(A,\ \mathbf{t}=A\mathbf{s}+\mathbf{e})$. \emph{Search-MLWE} is to recover $\mathbf{s}$; \emph{Decision-MLWE} is to distinguish $(A,\,A\mathbf{s}+\mathbf{e})$ from a uniform pair $(A,\,\mathbf{u})$.
\end{definition}

Two special cases bracket the definition: $k=1$ is RLWE (a single ring element), and $n=1$ with $R_q=\mathbb{Z}_q$ recovers plain LWE. Module-LWE interpolates between them, and it inherits a worst-case hardness guarantee of the same flavour as LWE: solving it on average is at least as hard as standard problems on \emph{module} lattices~\cite{langlois2015}. In Kyber the matrix is square ($m=k$), giving the compact form
\[
\mathbf{t} = A\mathbf{s} + \mathbf{e}, \qquad A\in R_q^{k\times k},\ \ \mathbf{s},\mathbf{e}\in R_q^{k}.
\]

\section{Why Module-LWE is considered a compromise}

Many people wrongly think that Module-LWE is simply "better" because it uses more structure. That is not the correct understanding. The real point is that Module-LWE reduces the strongest algebraic structure of pure RLWE while retaining the efficiency advantages. In other words, it trades off between:

\begin{itemize}
    \item the huge size of plain LWE,
    \item the strong structure of RLWE,
    \item and the practical middle ground provided by MLWE.
\end{itemize}

This balance is central to Kyber's design.

\section{Why the Kyber parameters are $k=2,3,4$}

The values $k=2,3,4$ correspond to different dimensions of the module, and are exactly the three security levels of Kyber~\cite{kyber2018,langlois2015}. In practice:

\begin{itemize}
    \item $k=2$ gives Kyber512,
    \item $k=3$ gives Kyber768,
    \item $k=4$ gives Kyber1024.
\end{itemize}

Increasing $k$ increases the size of the module and therefore increases the security level. The ring dimension and modulus stay fixed, but standardized ML-KEM parameter sets use different parameter tuples as well: in particular $\eta_1$ is $3,2,2$ and $(d_u,d_v)$ is $(10,4),(10,4),(11,5)$ for ML-KEM-512, -768, and -1024 respectively~\cite{fips203}.

\section{A complete toy MLWE experiment}

A useful exercise is to implement a toy version of MLWE in full.

\begin{enumerate}
    \item Define a ring $R_q$.
    \item Create a small polynomial matrix $A$.
    \item Create a polynomial vector $\mathbf{s}$.
    \item Create a noise vector $\mathbf{e}$.
    \item Compute $\mathbf{t}=A\mathbf{s}+\mathbf{e}$.
    \item Print the result and inspect its structure.
\end{enumerate}

This gives a practical understanding of why Kyber can be seen as a structured generalization of LWE. Written as pseudocode, using the same names $A$, $\mathbf{s}$, $\mathbf{e}$, $\mathbf{t}$ as the Sage code above, the same six steps are Algorithm~\ref{alg:mlwe-sample}.

\begin{algorithm}[H]
\caption{Module-LWE Sample Generation}
\label{alg:mlwe-sample}
\begin{algorithmic}[1]
\Require Ring $R_q = \mathbb{Z}_q[x]/(x^n+1)$; module dimension $k$
\Ensure MLWE sample $(A,\mathbf{t})$; secret $\mathbf{s}$
\State $A \gets$ uniformly random $k\times k$ matrix over $R_q$
\State $\mathbf{s} \gets$ length-$k$ vector of small (noise-sized) elements of $R_q$ \Comment{the secret}
\State $\mathbf{e} \gets$ length-$k$ vector of small (noise-sized) elements of $R_q$ \Comment{the error}
\State $\mathbf{t} \gets A\mathbf{s} + \mathbf{e}$ \Comment{a vector of $k$ ring elements, not one scalar}
\State \Return $(A,\mathbf{t})$ as the sample, $\mathbf{s}$ as the secret
\end{algorithmic}
\end{algorithm}

Real ML-KEM has this algorithmic shape with $n=256$ and $k=2,3,4$, but the standardized parameter sets differ in more than $k$; their sampling and compression parameters are specified in FIPS~203.

\section{Summary}

The key lesson is that Module-LWE is not merely "LWE with more polynomials". It is the middle ground between the two extremes:

\begin{table}[h]
\centering
\begin{tabular}{|c|c|c|}
\hline
Scheme & Secret object & Main idea \\
\hline
LWE & vector over $\mathbb{Z}_q$ & simple and general, but large \\
\hline
RLWE & single polynomial over $R_q$ & compact and fast, but highly structured \\
\hline
Module-LWE & vector of polynomials over $R_q$ & balanced compromise, used by Kyber \\
\hline
\end{tabular}
\end{table}

The important insight is that Kyber is not simply RLWE with a different name. It is the module version of the same noisy linear structure, arranged in a way that gives practical efficiency while preserving strong security assumptions.
